取 Schwarzschild 黑洞质量参数 \(M=GM_{\rm BH}/c^2>0\)，以下用几何单位 \(G=c=1\)，度规
\[
ds^2=-f(r)dt^2+f(r)^{-1}dr^2+r^2(d\theta^2+\sin^2\theta\,d\phi^2),\qquad f(r)=1-\frac{2M}{r}.
\]
只研究赤道面 \(\theta=\pi/2\)、外部区域 \(r>2M\)。单位静质量粒子的四速度按 \(u^\mu u_\mu=-1\) 归一，光线满足零测地线条件。无穷远静止观测者使用坐标时间 \(t\)。
(1) 从 Killing 守恒量 \(\mathcal E,\ell\) 推出有质量粒子的径向有效势，求所有定半径圆测地线的 \(\mathcal E^2,\ell^2\) 和存在范围，核对轨道角频率。
(2) 固定角动量作径向小扰动，分别求固有时和无穷远坐标时的径向本征频率；证明最内稳定圆轨道(ISCO)的半径，并求近圆轨道每一径向周期的近心点进动及弱场首项。
(3) 对零测地线独立推导光子球半径及其不稳定性，求来自无穷远的临界冲量参数与几何光学捕获截面，指出不能将有质量粒子的圆轨道条件直接用于光子。
(4) 计算圆轨道粒子相对当地静止观测者的切向速度、时间膨胀与单位静质量束缚能。若粒子在自身瞬时静止系沿顺/逆运动方向各发出频率 \(\nu_0\) 的切向光子，求无穷远观测频率 \(\nu_\infty^\pm\)；在 ISCO 给出两者的精确比值，并讨论远场极限与 \(r\downarrow3M\) 的适用性。
